%% paper83_beal_vermutung_de.tex
%% Die Beal-Vermutung und Verallgemeinerungen des Grossen Fermatschen Satzes (Deutsch)
%% Autor: Michael Fuhrmann
%% Specialist Mathematics Project, Build 166
%% Datum: 2026-03-12

\documentclass[12pt,a4paper]{amsart}
\usepackage{amsmath,amssymb,amsthm,mathtools,enumitem,booktabs,array,hyperref}
\usepackage[utf8]{inputenc}
\usepackage[T1]{fontenc}
\usepackage[ngerman]{babel}
\usepackage{geometry}
\geometry{margin=2.5cm}

%% Theorem-Umgebungen (Deutsch)
\newtheorem{theorem}{Satz}[section]
\newtheorem{satz}[theorem]{Satz}
\newtheorem{lemma}[theorem]{Lemma}
\newtheorem{korollar}[theorem]{Korollar}
\newtheorem{corollary}[theorem]{Korollar}
\newtheorem{proposition}[theorem]{Proposition}
\newtheorem{vermutung}[theorem]{Vermutung}
\newtheorem{conjecture}[theorem]{Vermutung}
\theoremstyle{definition}
\newtheorem{definition}[theorem]{Definition}
\newtheorem{definition_de}[theorem]{Definition}
\newtheorem{beispiel}[theorem]{Beispiel}
\theoremstyle{remark}
\newtheorem{bemerkung}[theorem]{Bemerkung}
\newtheorem{remark}[theorem]{Bemerkung}

\hypersetup{
  colorlinks=true,
  linkcolor=blue,
  citecolor=blue,
  urlcolor=blue
}

\begin{document}

\title{Die Beal-Vermutung\\
und Verallgemeinerungen des Großen Fermatschen Satzes:\\
Exponentielle diophantische Gleichungen\\
mit paarweisen Teilerfremdheitsbedingungen}

\author{Michael Fuhrmann}
\address{Specialist Mathematics Project, Build 166}
\email{specialist-maths@project.local}
\date{2026-03-12}

\begin{abstract}
Die Beal-Vermutung, unabhängig von Andrew Beal formuliert (1993) und verwandt
mit der Tijdeman--Zagier-Vermutung, behauptet: Falls $A^x + B^y = C^z$ mit
positiven ganzen Zahlen $A, B, C, x, y, z$ und $x, y, z \geq 3$, dann muss
$\gcd(A, B, C) > 1$ gelten, d.h.\ die drei Basen müssen einen gemeinsamen Teiler
größer als $1$ besitzen.

Wir entwickeln die Vermutung in ihrem historischen Kontext als Verallgemeinerung
des Großen Fermatschen Satzes, analysieren bekannte parametrische Lösungsfamilien
(alle mit $\gcd > 1$ oder einem Exponenten $< 3$) und präsentieren die wichtigsten
Teilresultate: Wiles' Beweis des GFS für $x = y = z$, den Darmon--Granville-Satz
über endlich viele primitive Lösungen bei festen Exponenten, die tiefe Verbindung
zur $abc$-Vermutung sowie die Ergebnisse von Poonen--Schaefer--Stoll.
Das Preisgeld von $1{.}000{.}000$ \$ von Andrew Beal wird ebenfalls besprochen.
\end{abstract}

\maketitle

\tableofcontents

%% =============================================================================
\section{Einleitung und geschichtlicher Hintergrund}
%% =============================================================================

Der Große Fermatsche Satz (GFS) besagt, dass $A^n + B^n = C^n$ für $n \geq 3$
keine Lösung in positiven ganzen Zahlen hat. Nach 358 Jahren wurde er 1995 von
Andrew Wiles~\cite{Wiles95} bewiesen. Die natürliche Frage lautet: Was geschieht,
wenn die Exponenten verschieden sein dürfen?

\begin{vermutung}[Beal-Vermutung, 1993]\label{verm:beal}
Falls $A, B, C, x, y, z$ positive ganze Zahlen mit $x, y, z \geq 3$ und
\[
  A^x + B^y = C^z,
\]
dann gilt $\gcd(A, B, C) > 1$.
\end{vermutung}

\begin{definition_de}
Eine Lösung $(A, B, C, x, y, z)$ von $A^x + B^y = C^z$ heißt \emph{primitiv}, falls
$\gcd(A, B, C) = 1$.
\end{definition_de}

Beals Vermutung behauptet: Es gibt keine primitiven Lösungen, wenn $x, y, z \geq 3$.

\begin{bemerkung}[Geschichte]
Andrew Beal ist Bankier und Hobbymathematiker, der die Vermutung 1993 nach
ausgedehnter Computersuche ohne Gegenbeispiel aufstellte. Das Problem wurde
vorher auch betrachtet:
\begin{itemize}
  \item Robert Tijdeman und Don Zagier formulierten dieselbe Vermutung; sie heißt
        manchmal auch \emph{Tijdeman--Zagier-Vermutung}.
  \item Der GFS ist der Spezialfall $x = y = z \geq 3$.
  \item Die Vermutung steht in der Tradition: Diophantische Gleichungen hohen
        Grades haben sehr wenige Lösungen.
\end{itemize}
Beal bietet ein Preisgeld für Beweis oder Widerlegung: derzeit $1{.}000{.}000$~\$.
\end{bemerkung}

%% =============================================================================
\section{Bezug zum Großen Fermatschen Satz}
%% =============================================================================

\begin{satz}[Fermat--Wiles, 1995]\label{satz:GFS}
Es gibt keine positiven ganzzahligen Lösungen von $A^n + B^n = C^n$ für $n \geq 3$.
\end{satz}

\begin{korollar}
Der Große Fermatsche Satz ist ein Spezialfall der Beal-Vermutung (der Fall $x = y = z$).
\end{korollar}

\begin{proof}
Für $x = y = z = n \geq 3$ besagt die Beal-Vermutung: $A^n + B^n = C^n$ impliziert
$\gcd(A,B,C) > 1$. Der GFS sagt, es gibt gar keine Lösungen, was stärker ist.
Also impliziert der GFS diesen Spezialfall.
\end{proof}

\begin{bemerkung}
Wiles' Beweis verwendet modulare elliptische Kurven und Galois-Darstellungen.
Dieselben Techniken lassen sich nicht direkt auf ungleiche Exponenten übertragen,
weshalb die Beal-Vermutung offen bleibt.
\end{bemerkung}

%% =============================================================================
\section{Bekannte parametrische Lösungen}
%% =============================================================================

Alle bekannten Lösungen von $A^x + B^y = C^z$ mit $x, y, z \geq 3$ haben
$\gcd(A,B,C) > 1$, was mit der Beal-Vermutung verträglich ist.

\begin{beispiel}[Standardparametrische Familien]\label{bsp:param}
\begin{enumerate}[label=(\alph*)]
  \item $3^3 + 6^3 = 3^5$: $27 + 216 = 243 = 3^5$. Hier $\gcd(3,6,3) = 3 > 1$.
  \item $2^6 + 2^6 = 2^7$: $64 + 64 = 128$, d.h.\ $(2^2)^3 + (2^2)^3 = 2^7$.
        $\gcd(4,4,2) = 2 > 1$.
\end{enumerate}
\end{beispiel}

\begin{satz}[Parametrische Lösungen mit gemeinsamen Teiler]\label{satz:param}
Für alle ganzen Zahlen $a, b$ mit $a^3 + b^3 = c^k$ für ganze $c, k \geq 3$
gilt $\gcd(a,b,c) > 1$, sofern die Lösung primitiv ist und der GFS angewandt wird.
Allgemeinere Familien: Für $d = \gcd(A,B,C) > 1$ schreibe $A = da$, $B = db$,
$C = dc$; dann wird $d^x a^x + d^y b^y = d^z c^z$, und für geeignete Wahl von
$a,b,c$ entstehen gültige Lösungen.
\end{satz}

%% =============================================================================
\section{Der Darmon--Granville-Satz}
%% =============================================================================

\begin{satz}[Darmon--Granville, 1995]\label{satz:DG}
Seien $p, q, r \geq 1$ ganze Zahlen mit $1/p + 1/q + 1/r < 1$. Dann hat die Gleichung
\[
  A^p + B^q = C^r
\]
nur endlich viele primitive Lösungen $(A, B, C) \in \mathbb{Z}^3$ mit $\gcd(A, B, C) = 1$.
\end{satz}

\begin{proof}[Beweisidee]
Darmon und Granville~\cite{DG95} nutzen Faltings' Theorem (frühere Mordell-Vermutung):
Die Gleichung $A^p + B^q = C^r$ definiert eine Kurve $\mathcal{C}$ von Geschlecht
$g \geq 2$, wenn $1/p + 1/q + 1/r < 1$ gilt (via Riemann--Hurwitz für eine
verzweigte Überlagerung der projektiven Geraden). Faltings' Satz liefert dann
nur endlich viele rationale Punkte.

Die Geschlechtsformel: Für $x^p + y^q = z^r$ hat ein glattes projektives Modell
Geschlecht $g \geq 2$, wenn $1/p + 1/q + 1/r < 1$.
\end{proof}

\begin{korollar}
Für $p, q, r \geq 3$: $1/p + 1/q + 1/r \leq 1$, mit Gleichheit nur für $(3,3,3)$.
Für $(3,3,3)$: GFS gilt. Für alle anderen $(p,q,r) \neq (3,3,3)$ mit $p,q,r \geq 3$:
Darmon--Granville liefert endlich viele primitive Lösungen.
\end{korollar}

\begin{bemerkung}
Darmon--Granville beweist \emph{Endlichkeit}, nicht \emph{Nichtexistenz}.
Die Beal-Vermutung behauptet, es gibt \emph{gar keine} primitiven Lösungen.
\end{bemerkung}

%% =============================================================================
\section{Fälle mit kleinen Exponenten}
%% =============================================================================

Wenn ein Exponent gleich $1$ oder $2$ ist, gibt es reichlich Lösungen.
Die Einschränkung $x,y,z \geq 3$ in der Beal-Vermutung ist wesentlich.

\begin{beispiel}[Exponent $= 2$]
$A^x + B^y = C^2$ mit $x,y \geq 2$: Es gibt unendlich viele primitive Lösungen.
Zum Beispiel: $2^3 + 1^2 = 3^2$ — check: $8 + 1 = 9$. Hier $\gcd(2,1,3) = 1$,
aber Exponent $2$ kommt vor, also kein Verstoß gegen Beal.
\end{beispiel}

\begin{satz}[GFS-Verbindungen für $z = 2$]\label{satz:z2}
Die Gleichung $A^p + B^q = C^2$ mit $p, q \geq 4$ hat keine primitiven Lösungen
für $p = q = 4$ (Euler--Fermat-Abstieg).
\end{satz}

%% =============================================================================
\section{Die \texorpdfstring{$abc$}{abc}-Vermutung und Beal}
%% =============================================================================

\begin{vermutung}[$abc$-Vermutung, Masser--Oesterlé 1985]\label{verm:abc}
Für jedes $\varepsilon > 0$ gibt es nur endlich viele Tripel $(a,b,c)$ paarweise
teilerfremder positiver ganzer Zahlen mit $a + b = c$ und
\[
  c > \mathrm{rad}(abc)^{1+\varepsilon},
\]
wobei $\mathrm{rad}(n) = \prod_{p \mid n} p$ das Produkt der verschiedenen Primteiler
von $n$ ist.
\end{vermutung}

\begin{bemerkung}
Die $abc$-Vermutung ist offen. Shinichi Mochizuki behauptete 2012 einen Beweis
(Inter-universelle Teichmüller-Theorie), der von der Gemeinschaft aber nicht
akzeptiert wurde; Scholze und Stix identifizierten 2018 eine Lücke.
\end{bemerkung}

\begin{satz}[$abc$ impliziert Beal]\label{satz:abcBeal}
Unter Annahme der $abc$-Vermutung: Die Beal-Vermutung gilt für alle bis auf endlich
viele primitive Tripel $(A, B, C)$ mit festen Exponenten $(x, y, z)$. Falls
$1/x + 1/y + 1/z < 1$, gibt es keine primitiven Lösungen für hinreichend große $A, B, C$.
\end{satz}

\begin{proof}
Sei $A^x + B^y = C^z$ mit $\gcd(A,B,C) = 1$. Setze $a = A^x$, $b = B^y$, $c = C^z$.
Dann $a + b = c$ und $\gcd(a,b,c) = 1$. Es gilt:
\[
  \mathrm{rad}(abc) = \mathrm{rad}(A^x B^y C^z) = \mathrm{rad}(ABC) \leq A \cdot B \cdot C
  \leq a^{1/x} \cdot b^{1/y} \cdot c^{1/z}.
\]
Die $abc$-Vermutung liefert $c < \mathrm{rad}(abc)^{1+\varepsilon} \leq
(a^{1/x} b^{1/y} c^{1/z})^{1+\varepsilon}$, also
\[
  c^{1-(1+\varepsilon)/z} < (ab)^{(1+\varepsilon)/\min(x,y)}.
\]
Für $1/x + 1/y + 1/z < 1$ und kleines $\varepsilon$ folgt eine polynomiale Schranke
für $c$, was nur endlich viele Lösungen erlaubt.
\end{proof}

%% =============================================================================
\section{Ergebnisse von Poonen--Schaefer--Stoll}
%% =============================================================================

\begin{satz}[Poonen--Schaefer--Stoll, 2007]\label{satz:PSS}
Die Gleichung $x^2 + y^3 = z^7$ hat genau $16$ primitive Lösungsfamilien (bis auf Vorzeichen).
Alle Lösungen erfüllen $\gcd(x,y,z) = 1$; da zwei Exponenten $< 3$ sind, widerspricht dies nicht der Beal-Vermutung.
\end{satz}

\begin{bemerkung}
Das Resultat aus~\cite{PSS07} verwendet Chabauty--Coleman-Methoden, Modulformen
und einen sorgfältigen Abstieg auf Jakobiannen von Kurven. Es zeigt, dass die
Theorie rationaler Punkte auf Kurven hohen Geschlechts (Faltings) effektiv
gemacht werden kann.
\end{bemerkung}

\begin{satz}[Bruin, 2003]\label{satz:bruin}
Die Gleichung $x^2 + y^4 = z^6$ hat keine primitiven Lösungen.
\end{satz}

%% =============================================================================
\section{Heuristische Dichte-Argumente}
%% =============================================================================

\begin{satz}[Dichte-Argument]\label{satz:dichte}
Heuristisch ist die Anzahl primitiver Lösungen von $A^x + B^y = C^z$ mit
$\max(A,B,C) \leq N$ von der Ordnung $O(N^{1/x + 1/y + 1/z - 1 + \varepsilon})$.
Für $1/x + 1/y + 1/z < 1$ geht diese Schätzung gegen $0$ für $N \to \infty$,
was die Beal-Vermutung heuristisch stützt.
\end{satz}

\begin{bemerkung}
Dieses Heuristikargument ist kein Beweis, erklärt aber, warum die Beal-Vermutung
plausibel ist: Die „erwartete" Anzahl von Lösungen sinkt mit wachsenden Exponenten.
\end{bemerkung}

%% =============================================================================
\section{Verbindungen zu Modularität und elliptischen Kurven}
%% =============================================================================

\begin{satz}[Ribet--Wiles für $A^p + B^p = C^2$]\label{satz:rpw}
Es gibt keine primitiven Lösungen von $A^p + B^p = C^2$ für Primzahlen $p \geq 7$
mit $2 \mid C$ und $4 \mid A$.
\end{satz}

\begin{satz}[Verallgemeinerter Fermat bei festen Exponenten]\label{satz:genFerm}
\begin{enumerate}[label=(\roman*)]
  \item $A^n + B^n = C^2$: Keine primitiven Lösungen für $n = 3$ (Euler) und
        $n = 4$ (Fermat-Abstieg).
  \item $A^3 + B^3 = C^n$: Keine primitiven Lösungen für $n = 3$ (Euler);
        für $n = 5, 7$ gelöst durch Bruin~\cite{Bruin03}.
\end{enumerate}
\end{satz}

\begin{bemerkung}[Frey-Kurven-Verbindung]
Für $A^p + B^q = C^r$ kann manchmal eine Frey--Hellegouarch-Kurve
$E\colon y^2 = x(x - A^p)(x + B^q)$ konstruiert werden.
Die Methode funktioniert sauber für $p = q = r$ (GFS), stößt aber bei ungleichen
Exponenten auf wesentliche Hindernisse.
\end{bemerkung}

%% =============================================================================
\section{Das Preisgeld und aktueller Stand}
%% =============================================================================

\begin{itemize}
  \item Andrew Beal kündigte die Vermutung 1993 nach ausgedehnter Computersuche
        ohne Gegenbeispiel an.
  \item Aktuelles Preisgeld: $1{.}000{.}000$~\$ (verwahrt von der American Mathematical Society).
  \item Das Preisgeld wird für einen begutachteten Beweis oder ein Gegenbeispiel ausgelobt.
  \item Computerverifikation: Keine primitiven Lösungen gefunden mit
        $\max(A,B,C) \leq 10^{12}$ und $3 \leq x,y,z \leq 12$.
\end{itemize}

\begin{vermutung}[Beal, aktueller Status]\label{verm:beal_status}
Die Beal-Vermutung (Vermutung~\ref{verm:beal}) ist offen für alle $x, y, z \geq 3$
über den positiven ganzen Zahlen. Keine primitiven Lösungen sind bekannt, und kein
Beweis der Nichtexistenz existiert.
\end{vermutung}

%% =============================================================================
\section{Verallgemeinerungen: Die verallgemeinerte Fermat-Gleichung}
%% =============================================================================

\begin{definition_de}
Die \emph{verallgemeinerte Fermat-Gleichung} (VFG) fragt nach ganzzahligen Lösungen von
\[
  x^p + y^q + z^r = 0
\]
mit $\gcd(x,y,z) = 1$ und $p, q, r \geq 2$.
\end{definition_de}

\begin{satz}[Beukers, 1998]\label{satz:beukers}
Die verallgemeinerte Fermat-Gleichung $x^2 + y^3 + z^5 = 0$ hat genau $27$
primitive ganzzahlige Lösungen (bis auf Gesamtvorzeichenwechsel).
\end{satz}

\begin{satz}[Klassifikation der VFG nach Euler-Charakteristik]\label{satz:euler}
Sei $\chi(p,q,r) = 1/p + 1/q + 1/r - 1$. Dann:
\begin{enumerate}[label=(\roman*)]
  \item $\chi > 0$ (sphärisch): $(2,2,r)$, $(2,3,3)$, $(2,3,4)$, $(2,3,5)$ —
        unendlich viele Lösungen.
  \item $\chi = 0$ (euklidisch): $(2,3,6)$, $(2,4,4)$, $(3,3,3)$ —
        endlich viele primitive Lösungen.
  \item $\chi < 0$ (hyperbolisch): Alle anderen $(p,q,r)$ mit $p,q,r \geq 2$ —
        endlich viele nach Faltings.
\end{enumerate}
\end{satz}

%% =============================================================================
\section{Zusammenfassung}
%% =============================================================================

Die Beal-Vermutung ist eine natürliche Verallgemeinerung des Großen Fermatschen Satzes
auf möglicherweise ungleiche Exponenten. Während Wiles' GFS-Beweis ein Meilenstein war,
lässt er die Beal-Vermutung vollständig offen.

Die Evidenz zugunsten der Vermutung ist stark:
\begin{itemize}
  \item Keine primitiven Lösungen mit $x,y,z \geq 3$ wurden je gefunden.
  \item Heuristische Argumente sagen ihre Nichtexistenz stark voraus.
  \item Die $abc$-Vermutung impliziert Nichtexistenz für große Lösungen.
  \item Darmon--Granville sichert Endlichkeit für jedes feste Tripel $(x,y,z)$.
\end{itemize}

Ein Beweis würde neue Methoden jenseits des Modularitäts-/Faltings-Werkzeugs erfordern,
vermutlich mit Galois-Darstellungen, $p$-adischer Hodge-Theorie und gänzlich neuen
Ideen über exponentielle diophantische Gleichungen.

\begin{thebibliography}{99}

\bibitem{Wiles95}
A.~Wiles,
\textit{Modular elliptic curves and Fermat's last theorem},
Ann.\ of Math.\ (2) \textbf{141} (1995), Nr.\ 3, 443--551.

\bibitem{DG95}
H.~Darmon, A.~Granville,
\textit{On the equations $z^m = F(x,y)$ and $Ax^p + By^q = Cz^r$},
Bull.\ London Math.\ Soc.\ \textbf{27} (1995), Nr.\ 6, 513--543.

\bibitem{TW95}
R.~Taylor, A.~Wiles,
\textit{Ring-theoretic properties of certain Hecke algebras},
Ann.\ of Math.\ (2) \textbf{141} (1995), Nr.\ 3, 553--572.

\bibitem{PSS07}
B.~Poonen, E.~F.\ Schaefer, M.~Stoll,
\textit{Twists of $X(7)$ and primitive solutions to $x^2 + y^3 = z^7$},
Duke Math.\ J.\ \textbf{137} (2007), Nr.\ 1, 103--158.

\bibitem{Bruin03}
N.~Bruin,
\textit{Chabauty methods and covering techniques applied to generalized Fermat equations},
CWI Tract \textbf{133}, Centrum voor Wiskunde en Informatica, Amsterdam, 2002.

\bibitem{Cohen07}
H.~Cohen,
\textit{Number Theory, Vol.\ II: Analytic and Modern Tools},
Graduate Texts in Mathematics \textbf{240}, Springer, 2007.

\bibitem{Masser85}
D.~W.\ Masser,
\textit{Open problems},
in \textit{Proc.\ Symp.\ Analytic Number Theory} (W.~W.~L.\ Chen, Hrsg.),
Imperial College, London, 1985.

\bibitem{Oesterle}
J.~Oesterlé,
\textit{Nouvelles approches du ``théorème'' de Fermat},
Séminaire Bourbaki 1987--88, Exposé \textbf{694}, Astérisque \textbf{161--162} (1988), 165--186.

\bibitem{Beal2016}
A.~Beal,
\textit{Beal Conjecture},
\url{http://www.bealconjecture.com}, Zugriff 2026.

\end{thebibliography}

\end{document}
